\section{Polynomial boundary calculus}\label{sec:boundary}

The determinant side of the companion problem has a uniform solution. We first expand all boundary functions in a fixed basis of theta minors, then invert the expansion and fold indices beyond its dimension. For integers $0\le j\le k$, write
\begin{equation}\label{eq:Adef}
A_{k,j}(q)=(-1)^jq^{\binom{k-j}2}\qbin{k+j}{2j}.
\end{equation}

\subsection{A uniform recurrence for boundary determinants}
Retain $d=m+1$, $L=2d+1$ and $w=q^{1/2}$. For $j\ge0$, put
\[
W_j(y)=\prod_{s=0}^{j-1}(1+wq^sy)(1+wq^s/y),\qquad \Psi^{[j]}_a(x)=\mu_x\left[\frac{y^a-y^{-a}}{W_j(y)}\right].
\]
The empty product gives $\Psi^{[0]}_a=\Phi_a$. Let $U_d(t)$ be the $d\times d$ tridiagonal matrix with off-diagonal entries $t$ and diagonal $1+t^2,\dots,1+t^2,1$. Define
\[
V_0=I_d,\qquad V_j=U_d(w)U_d(wq)\cdots U_d(wq^{j-1})\quad(j\ge1),
\]
where $I_d$ is the identity matrix.

\begin{lemma}[Boundary determinant expansion]\label{lem:bdexp}
For generic $x$ and every integer $j\ge0$,
\begin{equation}\label{eq:bdexp}
\HL_m(q^j;x;q)=\sum_{k=0}^m(-1)^k(V_j)_{d-k,d}\,J_k(x).
\end{equation}
In particular, for $m\ge2$,
\begin{equation}\label{eq:bdexp2}
\HL_m(q^2;x;q)-(1+q+q^2)\HL_m(q;x;q)+q\HL_m(1;x;q)=q^2J_2(x).
\end{equation}
\end{lemma}

\begin{proof}
In the bilateral formula used in the proof of \cref{lem:boundary}, take $t=q^j$. The finite Pochhammer quotient telescopes over the indices $-j,\dots,j-1$. Combined with $P_{q^j}/P_1$ it gives $1/W_j(y)$, where $y=xq^\ell$. Thus
\[
\HL_m(q^j;x;q)=P_1(x)\frac{\det(\Psi^{[j]}_a(x_i))_{i,a=1}^m}{\delta_C(x)}.
\]

Fix $j\ge1$ and set $t=wq^{j-1}$. Direct cancellation shows
\[
\Psi^{[j]}_{d+1}+t\Psi^{[j]}_d=\mu_x\left[\frac{y^{d+1}}{(1+ty)W_{j-1}(y)}-\frac{y^{-d-1}}{(1+t/y)W_{j-1}(y)}\right]=0.
\]
To verify the last equality, apply $\mu_x[f(y)]=q^{L/2}\mu_x[y^Lf(qy)]$ to the second term and use
\[
(1+t/(qy))\,W_{j-1}(qy)=\frac{(1+ty)\,W_{j-1}(y)}{wy}.
\]
Multiplication by the last factor of $W_j$ also gives
\[
\Psi^{[j-1]}_a=(1+t^2)\Psi^{[j]}_a+t\bigl(\Psi^{[j]}_{a-1}+\Psi^{[j]}_{a+1}\bigr),\qquad \Psi^{[j]}_0=0.
\]
The boundary relation therefore closes this recurrence at $a=d$:
\[
\bigl(\Psi^{[j]}_1,\dots,\Psi^{[j]}_d\bigr)^\T=V_j^{-1}(\Phi_1,\dots,\Phi_d)^\T.
\]
Each $U_d(t)$ has determinant $1$, and the complementary-minor formula and the Cauchy--Binet formula give \eqref{eq:bdexp}.

The last column of $V_1$ has entries $w,1$ in its last two rows. The last column of $V_2$ has entries
\[
q^2,\qquad w(1+q+q^2),\qquad 1+q^2
\]
in its last three rows and zeros elsewhere. Substitution in \eqref{eq:bdexp} proves \eqref{eq:bdexp2}.
\end{proof}

\subsection{The half-line kernel and Bailey inversion}

\begin{proposition}[Explicit boundary coefficients]\label{prop:coeffs}
Let $d=m+1$ and $0\le h\le j\le m$. Then
\begin{equation}\label{eq:coeffs}
(V_j)_{d-h,d}=q^{h^2/2}D_{j,h}(q),\qquad D_{j,h}(q)=\qbin{2j}{j-h}-q^{2h+1}\qbin{2j}{j-h-1}.
\end{equation}
Gaussian binomial coefficients outside their usual range are zero.
\end{proposition}

\begin{proof}
We first work on the half-line with coordinates $0,1,2,\dots$, where coordinate zero corresponds to the last row of the finite matrix. Let $\mathsf U(t)$ be tridiagonal, with off-diagonal entries $t$, diagonal entry $1$ at zero, and diagonal entries $1+t^2$ elsewhere. All matrix products below have finite bandwidth, so each entry is a finite sum. Put
\[
f_j(a)=q^{a^2/2}\qbin{2j}{j+a}\quad(a\in\Z),\qquad K_j(h,\ell)=f_j(h-\ell)-q^{h+1/2}f_j(h+\ell+1).
\]
The finite $q$-binomial theorem gives the Laurent polynomial identity
\[
\sum_{a\in\Z}f_j(a)\,y^a=\prod_{s=0}^{j-1}(1+wq^sy)(1+wq^s/y).
\]
Consequently, with $t_j=wq^j$,
\[
f_{j+1}(a)=(1+t_j^2)f_j(a)+t_j\bigl(f_j(a-1)+f_j(a+1)\bigr).
\]
The same recurrence holds for $K_j(h,\ell)$ in the second coordinate. The elementary ratio of consecutive Gaussian binomial coefficients gives
\[
(1-q^{j+h+1})f_j(h+1)=q^{h+1/2}(1-q^{j-h})f_j(h),
\]
also at the endpoints, where the appropriate terms vanish. Equivalently,
\[
K_j(h,-1)=-t_jK_j(h,0).
\]
This relation changes the diagonal coefficient at zero from $1+t_j^2$ to $1$. Since $K_0(h,\ell)=\delta_{h\ell}$ for $h,\ell\ge0$, induction proves
\[
K_j(h,\ell)=\bigl(\mathsf U(w)\,\mathsf U(wq)\cdots\mathsf U(wq^{j-1})\bigr)_{h,\ell}.
\]
For the column $\ell=0$ and $j\le m$, no matrix path reaches a coordinate above $m$, so this column agrees with that of the finite product $V_j$. Substituting $\ell=0$ gives \eqref{eq:coeffs}.
\end{proof}

\begin{lemma}[Uniform determinant inverse]\label{lem:inverse}
For generic $x$ and every $0\le k\le m$,
\begin{equation}\label{eq:inverse}
\sum_{j=0}^kA_{k,j}(q)\,\HL_m(q^j;x;q)=q^{k^2/2}J_k(x).
\end{equation}
\end{lemma}

\begin{proof}
By \cref{lem:bdexp} and \cref{prop:coeffs}, it suffices to show that
\begin{equation}\label{eq:inv-coeff}
\sum_{j=h}^kA_{k,j}(q)\,D_{j,h}(q)=(-1)^k\delta_{kh}\qquad(0\le h\le k).
\end{equation}
The case $h=k$ is immediate. Suppose $n=k-h\ge1$. Direct cancellation gives
\[
D_{j,h}(q)=\frac{(1-q^{2h+1})(q;q)_{2j}}{(q;q)_{j-h}(q;q)_{j+h+1}}.
\]
Define a polynomial of degree $n-1$ by
\[
P(T)=\prod_{a=2k-n+2}^{2k}(1-q^aT).
\]
With $\ell=k-j$, the left-hand side of \eqref{eq:inv-coeff} becomes
\[
\frac{(-1)^k(1-q^{2h+1})}{(q;q)_n}\sum_{\ell=0}^n(-1)^\ell q^{\binom\ell2}\qbin n\ell P(q^{-\ell}).
\]
For each monomial $T^a$ of $P$ one has $0\le a<n$, and the finite $q$-binomial theorem evaluates its contribution as $(q^{-a};q)_n=0$. This proves \eqref{eq:inv-coeff} and the lemma.
\end{proof}

The coefficient inversion \eqref{eq:inv-coeff} is a rescaled form of classical Bailey inversion. Indeed, put
\[
\beta_j=\frac{\HL_m(q^j;x;q)}{(q;q)_{2j}},\qquad \alpha_h=\frac{1-q^{2h+1}}{1-q}(-1)^hq^{h^2/2}J_h(x).
\]
For $0\le j\le m$, \cref{prop:coeffs} gives
\[
\beta_j=\sum_{h=0}^j\frac{\alpha_h}{(q;q)_{j-h}(q^2;q)_{j+h}},
\]
which is the defining Bailey relation relative to $q$; see \cite[Lemma~1]{MZ10} for the more general well-poised inversion. We include the direct proof above to make the normalization and the finite range explicit.

\subsection{The finite kernel and boundary powers beyond the rank}
The half-line kernel also gives an exact formula for the finite matrix with no restriction on the number of factors. This supplies boundary-shift relations even when the index exceeds the rank parameter.

\begin{proposition}[Finite boundary kernel]\label{prop:kernel}
Put $d=m+1$ and $L=2d+1$. For $j\ge0$ and $a\in\Z$ define
\[
f_j(a)=q^{a^2/2}\qbin{2j}{j+a},\qquad \gamma_i(h)=q^{i(h+1/2)+Li(i-1)/2}\quad(i\in\Z).
\]
For $0\le h,\ell<d$,
\begin{equation}\label{eq:kernel}
(V_j)_{d-h,d-\ell}=\sum_{i\in\Z}\gamma_i(h)\Bigl(f_j(h-\ell+iL)-q^{h+iL+1/2}f_j(h+\ell+1+iL)\Bigr).
\end{equation}
The sum is finite because $f_j(a)=0$ for $|a|>j$.
\end{proposition}

\begin{proof}
Write $K_j(a,\ell)=f_j(a-\ell)-q^{a+1/2}f_j(a+\ell+1)$, now allowing $a\in\Z$, and denote the right-hand side of \eqref{eq:kernel} by $\widehat K_j(h,\ell)$. The recurrence and the lower boundary calculation in the proof of \cref{prop:coeffs} remain valid for every integer $a$. Thus, with $t_j=wq^j$,
\[
\widehat K_j(h,-1)=-t_j\widehat K_j(h,0).
\]
At the upper boundary, the two arguments of $f_j$ in $K_j(h+iL,d)$ differ by $L$. Since
\[
\gamma_{i+1}(h)=q^{h+iL+1/2}\gamma_i(h),
\]
successive terms cancel, giving $\widehat K_j(h,d)=0$. The free three-term recurrence therefore becomes multiplication on the right by $U_d(t_j)$ in reversed coordinates: the lower boundary has diagonal entry $1$, and the upper boundary is truncated. Finally, $\widehat K_0(h,\ell)=\delta_{h\ell}$, since in the stated range only the term with $i=0$ and $h=\ell$ contributes. Induction on $j$ proves the formula.
\end{proof}

\begin{theorem}[Affine folding of the boundary inverse]\label{thm:folding}
For generic $x$ and every $k\ge0$, put
\[
\mathcal Q_{m,k}(x;q)=q^{-k^2/2}\sum_{j=0}^kA_{k,j}(q)\,\HL_m(q^j;x;q).
\]
Write $k=bL+s$ with $b\ge0$ and $0\le s<L$. Then
\begin{equation}\label{eq:folding}
\mathcal Q_{m,k}(x;q)=(-1)^bq^{b(s+1/2)+L\binom b2}\begin{cases}J_s(x),&0\le s<d,\\ 0,&s=d,\\ -J_{L-1-s}(x),&d<s<L.\end{cases}
\end{equation}
In particular,
\[
\mathcal Q_{m,m+1}=0,\qquad \mathcal Q_{m,k+L}=-q^{k+1/2}\mathcal Q_{m,k}\quad(k\ge0).
\]
\end{theorem}

\begin{proof}
For $a\ge0$, the coefficient calculation \eqref{eq:inv-coeff} gives
\[
\sum_{j=0}^kA_{k,j}(q)\,K_j(a,0)=(-1)^kq^{k^2/2}\delta_{ak};
\]
this includes $a>k$, since all summands then vanish. The symmetry $f_j(-a)=f_j(a)$ also gives
\[
K_j(-a-1,0)=-q^{-a-1/2}K_j(a,0).
\]
Substitute the finite kernel into \cref{lem:bdexp} and perform the sum over $j$. For fixed $h$, only image indices satisfying $h+iL=k$ or $h+iL=-k-1$ can survive.

If $s<d$, the first equation has the unique solution $h=s$, $i=b$, and the coefficient of $J_h$ is $(-1)^{k+h}\gamma_b(h)$; this is the first case of \eqref{eq:folding}. If $s>d$, the second equation has the unique solution $h=L-1-s$, $i=-b-1$. Its coefficient is $-(-1)^{k+h}\gamma_{-b-1}(h)q^{-k-1/2}$, which simplifies to the third case. For $s=d$, neither equation has a solution with $0\le h<d$. The stated quasi-periodicity follows directly.
\end{proof}

\begin{corollary}[A boundary-shift relation]\label{cor:bdrel}
Let $m\ge1$, $r=2m+1$ and $v\in\Q$. Then
\begin{equation}\label{eq:bdrel}
\sum_{j=0}^{m+1}A_{m+1,j}(q)\,(-q^{1/2-v};q)_j\,F_{r,B_v+je_r,0}(q)=0.
\end{equation}
\end{corollary}

\begin{proof}
By \cref{lem:boundary}, each term after the coefficient $A_{m+1,j}$ equals $(-q^{1/2-v};q)_\infty\HL_m(q^j;q^v\one;q)$. Apply \cref{thm:folding} with $k=m+1$.
\end{proof}

\subsection{A minimal basis for the boundary shifts}
The finite kernel gives reductions directly among boundary-shifted Nahm sums, without introducing companion sums. The coefficients are polynomials in $q$, despite the negative powers in the inverse formula.

\begin{theorem}[Boundary basis]\label{thm:basis}
For $j\ge0$ and $0\le\ell\le m$, define
\begin{equation}\label{eq:Cjl}
C_{j,\ell}(q)=\sum_{h=\ell}^m(-1)^hq^{-h^2/2}(V_j)_{d-h,d}\,A_{h,\ell}(q),\qquad d=m+1.
\end{equation}
Then $C_{j,\ell}\in\Z[q]$ and, for generic $x$,
\begin{equation}\label{eq:basis}
\HL_m(q^j;x;q)=\sum_{\ell=0}^mC_{j,\ell}(q)\,\HL_m(q^\ell;x;q).
\end{equation}
For each fixed $q$ with $0<|q|<1$, the $m+1$ functions $\HL_m(1;x;q),\dots,\HL_m(q^m;x;q)$ are linearly independent over $\C$ as meromorphic functions of $x$. Consequently these reductions are unique, and no nonzero relation with coefficients independent of $x$ involves only the first $m+1$ boundary shifts. The relation at $j=m+1$ is, up to a scalar factor,
\[
\sum_{j=0}^{m+1}A_{m+1,j}(q)\,\HL_m(q^j;x;q)=0.
\]
\end{theorem}

\begin{proof}
Substitute \cref{lem:inverse} into \eqref{eq:bdexp}; interchanging two finite sums gives \eqref{eq:Cjl} and \eqref{eq:basis}.

A matrix path contributing to $(V_j)_{d-h,d}$ has at least $h$ off-diagonal steps. Their time indices are distinct and nonnegative, and a step at time $s$ has weight $w^{2s+1}$, so their total exponent is at least $1+3+\cdots+(2h-1)=h^2$. Every diagonal contribution has even exponent in $w$, while the number of off-diagonal steps has the parity of $h$. Thus
\[
q^{-h^2/2}(V_j)_{d-h,d}\in\Z[q].
\]
Since $A_{h,\ell}\in\Z[q]$, the same holds for $C_{j,\ell}$.

To prove independence, fix $q$ and recall that
\[
\Phi_a(x)=\sum_{n\in\Z}q^{Ln^2/2}\bigl(q^{an}x^{Ln+a}-q^{-an}x^{Ln-a}\bigr),\qquad 1\le a\le d,\quad L=2d+1.
\]
The residue classes $\pm1,\dots,\pm d$ modulo $L$ are distinct, so the coefficient of $x^a$ in a linear combination of these Laurent series is exactly the coefficient of $\Phi_a$. Thus $\Phi_1,\dots,\Phi_d$ are linearly independent on $\C^\times$.

There exist $d$ evaluation points at which the matrix of values of $\Phi_1,\dots,\Phi_d$ is invertible. Evaluating the $d$ maximal minors on the $d$ subsets of $d-1=m$ of these points produces, up to signs, the cofactor matrix, which is also invertible. Hence these minors are linearly independent as functions of $m$ variables. Multiplication by the common nonzero meromorphic factor $P_1(x)/\delta_C(x)$ proves the independence of $J_0,\dots,J_m$.

For $0\le j\le m$, the coefficient matrix in \eqref{eq:bdexp} is triangular: the coefficient of $J_h$ vanishes for $h>j$, and the coefficient of $J_j$ is $(-1)^jq^{j^2/2}\ne0$. Therefore the first $m+1$ boundary functions are independent. The remaining assertions follow from this independence and \cref{thm:folding} with $k=m+1$.
\end{proof}

\begin{corollary}[Polynomial reduction of every boundary shift]\label{cor:reduction}
Let $m\ge1$, $r=2m+1$, $v\in\Q$, $b=q^{1/2-v}$ and $F_j=F_{r,B_v+je_r,0}$. For every integer $j\ge0$,
\begin{equation}\label{eq:reduction}
(-b;q)_jF_j=\sum_{\ell=0}^mC_{j,\ell}(q)\,(-b;q)_\ell F_\ell.
\end{equation}
In particular, $C_{j,\ell}=\delta_{j\ell}$ for $0\le j\le m$, and $C_{m+1,\ell}=(-1)^mA_{m+1,\ell}$.
\end{corollary}

\begin{proof}
\Cref{lem:boundary} gives $(-b;q)_jF_j=(-b;q)_\infty\HL_m(q^j;q^v\one;q)$. Apply \eqref{eq:basis} and take the removable limit at coincident variables. The last assertions follow from the uniqueness in \cref{thm:basis}.
\end{proof}

The minimality statement concerns free variables $x$. We do not claim independence after setting $x=q^v\one$, or after fixing one of the modular sectors.

\subsection{A polynomial recurrence}
For indeterminates $e,c$, define
\begin{equation}\label{eq:Rdef}
R_k(e;c)=\sum_{j=0}^kA_{k,j}(q)\,(-c;q)_j\,e^j.
\end{equation}
The specialization $c=w/u$ occurs in the companion certificates. The following recurrence holds before that specialization.

\begin{proposition}[Boundary polynomial recurrence]\label{prop:recurrence}
One has $R_0(e;c)=1$, $R_1(e;c)=1-(1+c)e$, and, for $k\ge2$,
\begin{multline}\label{eq:recurrence}
R_k(e;c)=(1+q)q^{k-1}R_{k-1}(e/q;c)-q^{2k-2}R_{k-2}(e/q^2;c)\\
-(1+c)q^{2k-2}e\,R_{k-1}(e/q^2;cq).
\end{multline}
The equality is an identity in $\Z[q,q^{-1},c,e]$, and its left-hand side belongs to $\Z[q,c,e]$.
\end{proposition}

\begin{proof}
Put $S_k(e)=\sum_{j=0}^kA_{k,j}(q)e^j$ and define
\[
\mathcal F_n(X,s)=\sum_{a=0}^{\lfloor n/2\rfloor}q^{\binom a2}\qbin{n-a}{a}s^aX^{n-2a}.
\]
These are Cigler's $q$-Fibonacci polynomials, in the normalization $\mathcal F_n(X,s)=F_{n+1}(X,s/q)$ of \cite[Theorem~2.1]{Cigler03}. The $q$-Pascal identity gives
\[
\mathcal F_n(X,s)=X\mathcal F_{n-1}(X,qs)+s\mathcal F_{n-2}(X,qs)\qquad(n\ge2).
\]
Eliminating the odd index by two applications of this recurrence yields
\[
\mathcal F_{2k}(X,s)=X^2\mathcal F_{2k-2}(X,q^2s)+(1+q)s\,\mathcal F_{2k-2}(X,qs)-q^2s^2\mathcal F_{2k-4}(X,q^2s).
\]
When $X^2=-e$, the change of index $a=k-j$ gives $\mathcal F_{2k}(X,1)=S_k(e)$, and homogeneity gives $\mathcal F_{2h}(X,q^a)=q^{ah}S_h(e/q^a)$. Consequently,
\[
S_k(e)=(1+q)q^{k-1}S_{k-1}(e/q)-q^{2k-2}S_{k-2}(e/q^2)-q^{2k-2}e\,S_{k-1}(e/q^2).
\]
Let $\mathcal L_c$ be the $\Z[q,q^{-1}]$-linear map determined by $\mathcal L_c(e^j)=(-c;q)_je^j$. It commutes with rescaling of $e$ and satisfies, for every polynomial $P(e)$,
\[
\mathcal L_c(eP(e))=(1+c)\,e\,\mathcal L_{cq}(P(e)).
\]
Applying this map to the recurrence for $S_k$ proves \eqref{eq:recurrence}. The initial values follow from \eqref{eq:Rdef}.
\end{proof}

The recurrence is an exact relation among the coefficient polynomials. It does not by itself prove a recurrence among the Nahm sums to which those polynomials are applied.
